% ===========================================================================
%  sviluppatore/07-marketplace/06-codice.tex
% ===========================================================================
\section{Il codice che arriva da un altro}\label{sez:codice-marketplace}

Un acquisto installa codice Python scritto da un altro, sul computer dove
stanno le credenziali del broker, il token del sito e le strategie della
persona. La regola, in una riga: \textbf{il codice del marketplace fa conti
sui dati che il sistema gli passa}. Niente file, rete, programmi,
introspezione dell'interprete, codice costruito a runtime, modifiche al
sistema. La fanno rispettare tre strati indipendenti, perché nessuno basta da
solo.

\figuraTikz{codice-marketplace}
  {I tre strati. L'analisi gira due volte, alla pubblicazione e
   sul computer di chi compra; la verifica a ogni caricamento; la guardia
   per tutto il tempo in cui il codice gira.}
  {codice-marketplace}

\subsection{Primo strato: l'analisi statica}

\file{core/analisi\_codice.py} gira alla pubblicazione (rotta, pagina
\voce{Sell}, pubblicazione automatica: chi vende vede subito cosa non va) e,
soprattutto, \textbf{all'installazione sul computer di chi compra}, prima di
scrivere un solo file. Chi compra non dipende dal fatto che il sito abbia
controllato: il segreto con cui l'app parla al sito è nell'app distribuita, e
chiunque potrebbe pubblicare scavalcandola.

\begin{description}[leftmargin = 0pt, style = nextline, itemsep = 6pt]
\item[Gli import]
  Una lista di moduli \emph{ammessi}, non di vietati: la libreria standard di
  calcolo (\pkg{math}, \pkg{statistics}, \pkg{datetime}, \pkg{collections},
  \pkg{json}, \pkg{re}, \pkg{logging} senza \pkg{handlers} e \pkg{config},
  \ldots), \pkg{numpy}, \pkg{pandas}, \pkg{scipy}, \pkg{sklearn},
  \pkg{statsmodels}, \pkg{lightgbm} meno i sottomoduli che fanno I/O o
  eseguono codice (\cod{SOTTOMODULI\_VIETATI}), e delle parti di Lyra solo
  l'\term{SDK}: \cod{core.types}, \cod{core.indicators}, \cod{core.portfolio},
  \cod{core.volume\_profile}, \cod{core.exit\_policy}, \cod{strategies.base},
  \cod{strategies.descriptor}, \cod{approaches.descriptor},
  \cod{data.universe}, quattro funzioni di \cod{data.loader} (senza
  \cod{cache\_dir=}), \cod{data\_series} e \cod{custom\_series} di
  \cod{core.node\_data}, \cod{features.feature\_engine},
  \cod{features.panel\_features}, \cod{reporting.metrics}, tre funzioni di
  \cod{regime.logics.\_base}. Gli import relativi solo dentro la cartella
  dell'artefatto; mai un altro artefatto, mai sé stesso per nome.
\item[Le catene risolte]
  \cod{pd.core.config\_init.os.system} passa per nomi innocui e arriva a
  \pkg{os}. Ogni catena che parte da un modulo importato si segue sugli
  oggetti \emph{veri} (l'analisi importa le librerie ammesse): ogni modulo
  raggiunto deve essere ammesso, l'oggetto finale non deve venire da un modulo
  vietato (\cod{oggetto\_permesso}).
\item[I nomi]
  Builtin pericolosi (\cod{eval}, \cod{exec}, \cod{open}, \cod{globals},
  \ldots); dunder fuori da un elenco di protocolli; attributi che fanno I/O o
  portano all'interprete (\cod{read\_*}, \cod{to\_csv}, \cod{tb\_frame},
  \cod{f\_globals}, \cod{as\_strided}); \cod{getattr} con un nome calcolato
  (al suo posto \cod{core.types.by\_name}, che fa gli stessi controlli al
  momento della chiamata); stringhe di formato che attraversano attributi
  riservati (\cod{"\{0.\_\_globals\_\_\}"}); nomi di metodi passati come
  stringa (\cod{df.agg("to\_csv")}); assegnazioni e metodi che cambiano oggetti
  condivisi di Lyra (\cod{UNIVERSE.clear()}, \cod{np.mean = \ldots}).
\end{description}

\noindent
È volutamente severa: un falso allarme si risolve riscrivendo una riga, un via
libera sbagliato no. Ogni rifiuto dice file, riga e cosa fare invece.
\file{tests/campioni\_marketplace.py} raccoglie le vie di fuga note (\num{89})
e gli esempi di calcolo che devono passare; undici artefatti di calcolo della
libreria passano.

\begin{attenzione}{Le parole generiche si controllano solo sui moduli}
\cod{load}, \cod{config}, \cod{style}, \cod{options} sono pericolosi come
attributi di un modulo (\cod{np.load}, \cod{logging.config},
\cod{pd.options}) e parole qualunque come attributi propri
(\cod{self.config}, \cod{self.style}). Per questo stanno in
\cod{ATTRIBUTI\_DI\_MODULO} e non in \cod{ATTRIBUTI\_VIETATI}: chi aggiunge un
nome deve scegliere l'elenco giusto, o l'analisi rifiuterà mezza libreria.
\end{attenzione}

\subsection{Secondo strato: provenienza e verifica al caricamento}

L'installatore (\file{dashboard/marketplace\_install.py}) è l'unico punto
d'ingresso: pagina dell'artefatto, client del marketplace
(\rotta{POST}{/api/market/install/<id>}, che prima scriveva i file in una
cartella il cui nome veniva dalla richiesta), installazione per
l'amministratore. Oltre ad analizzare:

\begin{itemize}
  \item accetta nomi e percorsi \textbf{solo nelle forme previste}:
        identificatori, \file{<cartella>/<artefatto>/<file>.py} o
        \file{<cartella>/<artefatto>.py}; mai un modulo del sistema
        (\file{strategies/registry.py}, \file{weights/\_\_init\_\_.py}) né un
        \file{\_\_init\_\_.py} del pacchetto --- lo scrive Lyra, vuoto. Prima
        un progetto poteva scrivere \file{reporting/<qualsiasi>.py}, e da
        sorgente sovrascrivere il codice di Lyra;
  \item scrive il \textbf{proprio} \file{run\_backtest.py}
        (\cod{provenienza.lanciatore}): quello del pacchetto girerebbe come
        programma principale del backtest. Chi lancia il backtest lo riscrive
        se è stato cambiato, e non avvia mai un file del venditore come
        script; i registri fissano \cod{backtest\_script} e
        \cod{output\_path\_pattern} alla cartella dell'artefatto;
  \item in un progetto lascia \textbf{intatti} i componenti che la persona ha
        già, senza mescolarvi file del venditore;
  \item compila subito i \file{.pyc} (dopo, scritti dal codice del
        marketplace, la guardia li vieterebbe: un \file{.pyc} scritto da lui
        girerebbe al posto del sorgente controllato);
  \item annota ogni file in \file{data/marketplace/installati.json}
        (\file{core/provenienza.py}): impronta, artefatto, versione, versione
        delle regole.
\end{itemize}

\noindent
Ogni volta che uno di quei file sta per essere eseguito --- evento
\cod{exec} dell'audit hook, qualunque sia la strada dell'import ---
\cod{provenienza.verifica} confronta l'impronta: se il file o le regole sono
cambiati lo si ricontrolla; una modifica che resta nelle regole aggiorna il
registro, una che ne esce non gira (\cod{ImportError} con il motivo).

\subsection{Terzo strato: la guardia}

\file{core/guardia.py} è un audit hook (\simbolo{PEP~578}). Quando nel filo
che chiede un'operazione sensibile c'è un frame del codice del marketplace ---
riconosciuto dal \cod{\_\_file\_\_} dei globali del frame, anche sotto
funzioni di Python, di una libreria o di Lyra che lui ha chiamato --- valgono
le sue regole:

\begin{center}
\begin{tabular}{@{}L{36mm}L{102mm}@{}}
\toprule
\textbf{Che cosa} & \textbf{Regola} \\
\midrule
\rigaTenue
Rete & mai, nemmeno la risoluzione dei nomi \\
Processi, ambiente & mai \\
\rigaTenue
Disco & scritture e modifiche mai; letture solo dentro l'app, lo spazio di
  lavoro, Python e le cartelle di sistema, e mai \file{.env}, token,
  credenziali del broker, stato del live, runner, log \\
Interprete & mai pickle, \cod{gc.get\_objects}, codice o funzioni costruiti a
  runtime, \cod{settrace} \\
\rigaTenue
ctypes, frame, marshal & mai \emph{chiesti da lui}: threadpoolctl dentro
  scikit-learn li usa per sé, e funziona \\
\bottomrule
\end{tabular}
\end{center}

\noindent
Ogni blocco alza \cod{PermissionError} con il nome dell'artefatto e finisce in
\cod{ultimi\_blocchi()}, visibile in \rotta{GET}{/api/marketplace/diagnostica}.
L'hook si accende solo se nello spazio di lavoro c'è codice del marketplace
(gli hook non si tolgono, e senza nulla da sorvegliare non serve), e si
installa prima di qualunque artefatto: dall'avvio dell'app, dai pacchetti
degli artefatti, dal lanciatore, da \cod{run\_walkforward} e
\cod{run\_approach}. Su un carico pandas pesante il costo è nel rumore della
misura.

\begin{nota}{Il resto del confine}
\begin{itemize}
  \item Il backtest di un artefatto del marketplace parte con un ambiente
        senza chiavi, token e password (\cod{guardia.ambiente\_sottoprocesso});
        i segreti della sessione (\cod{LYRA\_TOKEN}, \cod{FLASK\_SECRET\_KEY})
        non arrivano a nessun sottoprocesso.
  \item I motori gli passano il portafoglio in sola lettura e gli indicatori
        in copia (\vediSez{live}).
  \item Metriche e indicatori arrivano come testo: se il testo è quello di
        un file del marketplace si compila con il suo percorso e la guardia
        li riconosce.
  \item Ciò che è stato comprato non si ripubblica: fuori dall'elenco di
        \voce{Sell}, rifiutato dalla rotta (anche dentro un progetto) e dalla
        pubblicazione automatica.
  \item \voce{Studio › Import\ldots}, codice portato dentro a mano: la stessa
        analisi produce un avviso, non un blocco. È una scelta di chi importa.
\end{itemize}
\end{nota}

\begin{attenzione}{Che cosa resta fuori}
Uno strato dentro il processo non ferma un bug di memoria nativo di
\pkg{numpy} o \pkg{pandas} sfruttato apposta: per questo l'analisi vieta le
primitive di memoria grezza e il processo del backtest non ha segreti
nell'ambiente. Le letture di \cod{os.environ} non passano per l'audit hook:
le blocca solo l'analisi statica. Chi aggiunge una funzione all'SDK deve
chiedersi se prende un percorso, un indirizzo o un oggetto con dentro il
broker: se sì, non va nell'SDK.
\end{attenzione}
